\documentclass[11pt,letterpaper]{article}
\usepackage[T1]{fontenc}
\usepackage{lmodern}
\usepackage[margin=1in]{geometry}
\usepackage{amsmath,amssymb,amsthm,mathtools}
\usepackage{enumitem}
\usepackage{microtype}
\usepackage{tikz}
\usepackage[colorlinks=true,linkcolor=blue,citecolor=blue,urlcolor=blue]{hyperref}
\usepackage{bookmark}
\pdftrailerid{}
\pdfsuppressptexinfo=15

\newtheorem{theorem}{Theorem}[section]
\newtheorem{proposition}[theorem]{Proposition}
\newtheorem{lemma}[theorem]{Lemma}
\newtheorem{corollary}[theorem]{Corollary}
\theoremstyle{definition}
\newtheorem{definition}[theorem]{Definition}
\theoremstyle{remark}
\newtheorem{remark}[theorem]{Remark}
\numberwithin{equation}{section}
\setlist{itemsep=3pt,topsep=5pt}
\hypersetup{pdftitle={Talagrand's discrete-convexity conjecture},
  pdfsubject={A positive solution to Talagrand's Research Problem 13.3.2}}

\newcommand{\unverified}[1]{\par\smallskip\noindent\fbox{\parbox{0.95\linewidth}{\textbf{Unverified step.} #1}}\par\smallskip}
\newcommand{\step}[1]{\par\smallskip\noindent\textbf{#1}\enspace}

\title{Talagrand's discrete-convexity conjecture}
\author{Rewritten for readability from the OpenAI manuscript of the same title}
\date{September 23, 2026}

\AddToHook{cmd/thebibliography/after}{\addcontentsline{toc}{section}{\refname}}

\begin{document}
\maketitle
\begin{abstract}
We prove Talagrand's discrete-convexity conjecture.
There is a universal integer $k$ such that, whenever an arbitrary family has
Bernoulli product measure at least $1-1/k$, the sets not contained
in a union of $k$ members admit a cover of total cost at most $1/2$
at the same density.
\end{abstract}

\section{Introduction}
\label{sec:introduction}

\paragraph{Setting.}
Let $[N]=\{1,\ldots,N\}$, and let $\mu_p$ be the probability measure on
$2^{[N]}$ under which each element is present independently with
probability $p\in(0,1)$. For $\mathcal D\subseteq2^{[N]}$ and an integer
$k\ge1$, the \emph{exceptional family} is
\begin{equation}
\label{eq:exceptional-definition}
E_k(\mathcal D)
=\bigl\{S\subseteq[N]:\text{no }D_1,\ldots,D_k\in\mathcal D
\text{ satisfy }S\subseteq D_1\cup\cdots\cup D_k\bigr\}.
\end{equation}
The $D_j$ may repeat and need not be disjoint. Two immediate facts:
$E_k(\mathcal D)$ is closed under supersets, and
$\varnothing\notin E_k(\mathcal D)$ when $\mathcal D\ne\varnothing$.

A family $\mathcal G$ of sets \emph{covers} a family $\mathcal A$ if every
$S\in\mathcal A$ contains some $I\in\mathcal G$. For $\rho\in(0,1)$, the
\emph{cost} of $\mathcal G$ at density $\rho$ is
$c_\rho(\mathcal G)=\sum_{I\in\mathcal G}\rho^{|I|}$, and $\mathcal A$ is
\emph{$\rho$-small} if it has a cover $\mathcal G$ with
$c_\rho(\mathcal G)\le1/2$. Since $\mu_\rho\{S:I\subseteq S\}=\rho^{|I|}$,
a union bound shows that a $\rho$-small family has $\mu_\rho$-measure at
most $1/2$; the cover is an explicit certificate of this.

\begin{theorem}
\label{thm:main}
Set $k=2^{75}$. For every integer $N\ge1$, every $p\in(0,1)$, and every
family $\mathcal D\subseteq2^{[N]}$,
\[
\mu_p(\mathcal D)\ge1-\frac1k
\quad\Longrightarrow\quad
E_k(\mathcal D)\text{ is }p\text{-small}.
\]
\end{theorem}

The constant is independent of $N$, $p$ and $\mathcal D$, and
$\mathcal D$ is arbitrary (no monotonicity). This is Talagrand's
conjecture \cite[Conjecture~7.1]{Talagrand10}, restated as Research
Problem~13.3.2 in~\cite{Talagrand21}, whose Proposition~13.3.3 explains
its connection to positive selector processes. The name reflects the
analogy with his Gaussian convexity question
\cite[Conjecture~2.1]{Talagrand10}, where sums of a balanced set play the
role of unions.

\paragraph{Proof at a glance.}
The logical order is
Lemma~\ref{lem:covering} $\Rightarrow$
Proposition~\ref{prop:half-density} $\Rightarrow$
Theorem~\ref{thm:main}.
\begin{itemize}
\item \emph{Weights (Section~\ref{sec:weights}, Steps 1--3).}
Let $\mu_q(\mathcal F)\ge1/2$. The biased Fourier expansion of
$\mathbf 1_{\mathcal F}$ yields weights $w(U)\ge0$, $U\subseteq[N]$, with
a global bound $\sum_Uq^{|U|}w(U)\le1$ (Parseval) and, for every
$S\in E_{32}(\mathcal F)$, a local bound
$\sum_{\varnothing\ne U\subseteq S}w(U)^{16}\ge2^{-32}$. The local bound
comes from a signed product measure on $32\times N$ arrays that vanishes
unless the $32$ rows jointly contain $S$.
\item \emph{From weights to a cover (Section~\ref{sec:weights}, Steps 4--7).}
Sets with $w(U)\ge2^{-64|U|}$ are used as generators; the global bound
makes them cheap at density $q/2^{70}$. The remaining sets are sorted
into bins by size $r$ and weight level $w(U)\approx2^{-h}$. If $S$
contains fewer than $2^{12h}$ sets from each bin, then each bin
contributes in total less than $2^{12h}\cdot2^{-16h}=2^{-4h}$ to the
local sum, which then falls below $2^{-32}$. So an exceptional set that
avoids the first generators contains at least $2^{12h}$ sets of some bin.
\item \emph{Covering lemma (Section~\ref{sec:covering}).}
Lemma~\ref{lem:covering} covers, cheaply and by a random construction,
all sets containing $2^{12h}$ members of one bin. This yields
Proposition~\ref{prop:half-density}: $E_{32}(\mathcal F)$ is
$(q/2^{70})$-small.
\item \emph{Density recovery (Section~\ref{sec:density}).}
With $L=2^{70}$, couple $L$ random sets, each of law $\mu_p$, whose union
has law $\mu_{Lp}$. The family $\mathcal F$ of unions of $L$ members of
$\mathcal D$ then has $\mu_{Lp}(\mathcal F)\ge1/2$ and
$E_{32}(\mathcal F)=E_{32L}(\mathcal D)$, so the proposition with $q=Lp$
gives density $q/L=p$.
\end{itemize}

The rest of this section is not used in the proof of
Theorem~\ref{thm:main}.

\paragraph{A consequence: positive selectors.}
(The source's family $T$ is renamed $\mathcal T$; $T$ is a number in
Section~\ref{sec:covering}.) Following \cite[Section~7]{Talagrand10}, let
$\mathcal T\subseteq[0,\infty)^N$ be nonempty,
\[
\phi(S)=\sup_{\tau\in\mathcal T}\sum_{i\in S}\tau_i,
\qquad M=\mathbb E_{\mu_p}\phi,
\]
and suppose $0<M<\infty$. Then $\{S:\phi(S)\ge k^2M\}$ is $p$-small for
$k=2^{75}$; the source notes that this recovers the conclusion of Park
and Pham \cite[Theorem~1.2]{ParkPham24}. Proof: put
$\mathcal D=\{D:\phi(D)<kM\}$. Markov's inequality gives
$\mu_p(\mathcal D)\ge1-1/k$. Since the $\tau_i$ are nonnegative, $\phi$
is monotone and $\phi(A\cup B)\le\phi(A)+\phi(B)$ (bound each $\tau$
separately, then take suprema). Hence
$S\subseteq D_1\cup\cdots\cup D_k$ with all $D_j\in\mathcal D$ forces
$\phi(S)<k\cdot kM$, so $\{\phi\ge k^2M\}\subseteq E_k(\mathcal D)$, and
the cover from Theorem~\ref{thm:main} covers it. This gives a covering
certificate, not just the measure bound that Markov's inequality provides.

\paragraph{Related work.}
(Attributions as given in the source.) Hua, Song and Tudose resolve the
Gaussian problem and derive a discrete statement with three unions and
covering density $p^C$ \cite[Corollaries~1.3--1.4]{HuaSongTudose26}.
Ascoli, He, Park and Talagrand give a threshold reformulation and prove
graph-containment cases
\cite[Proposition~1.5 and Theorems~1.10--1.13]{AscoliEtAl26}. Talagrand
introduced a fractional notion of smallness
\cite[Sections~6 and~7]{Talagrand10}, defined below; Li resolves the
fractional problem \cite[Theorem~1.9]{Li26}, and Park derives an integral
cover of $E_2(\mathcal D)$ at density $p/(C\max\{1,\log\log(1/p)\})$ when
$\mu_p(\mathcal D)>1/2$ \cite[Corollary~3.5]{Park26}. On methods: with
two rows, the signed measure of Section~\ref{sec:weights} is, per the
source, Li's signed kernel multiplied by the Bernoulli masses
\cite[Section~2]{Li26}; Park notes the extension to an even number of
unions \cite[Corollary~3.5, footnote~3]{Park26}; signed product operators
and biased Fourier analysis go back to \cite[Section~2.2]{Friedgut08} and
\cite[Sections~2.2 and~4.1]{EllisFilmusFriedgut12}; random unions also
appear in the covering constructions of
\cite[Sections~1.1 and~2]{FischerPerson25}.

\paragraph{A two-union corollary.}
This corollary trades fewer unions for a constant density loss and rests
on two cited results. Call $\mathcal A$ \emph{fractionally $\rho$-small}
if there are coefficients $a_I\ge0$ $(I\subseteq[N])$ with
\[
\sum_{I\subseteq S}a_I\ge1\quad(S\in\mathcal A),
\qquad
\sum_{I\subseteq[N]}\rho^{|I|}a_I\le\frac12 .
\]
(The source displays this at $\rho=p$ and uses it at other densities.)

\begin{corollary}[Two unions at constant density loss]
\label{cor:two-union}
There is a universal constant $C\ge1$ such that, for every integer $N\ge1$,
every $p\in(0,1)$, and every family $\mathcal D\subseteq2^{[N]}$,
\[
\mu_p(\mathcal D)>\frac12
\quad\Longrightarrow\quad
E_2(\mathcal D)\text{ is }(p/C)\text{-small}.
\]
\end{corollary}

\begin{proof}
We write $\sigma$ for the density the source calls $r$.
\begin{enumerate}
\item If $E_2(\mathcal D)=\varnothing$, the empty cover suffices.
Otherwise $E_2(\mathcal D)$ is nonempty, increasing, and proper: it
omits $\varnothing$, since $\mu_p(\mathcal D)>1/2$ gives
$\mathcal D\ne\varnothing$.
\item \emph{Cited.} By Li's result \cite[Corollary~3.1]{Li26} in Park's
fractional-cover formulation \cite[Theorem~3.4]{Park26},
$E_2(\mathcal D)$ is fractionally $\sigma$-small for $\sigma=p/2$.
\item Replace each $a_I$ by $\min\{a_I,1\}$. The cost does not increase.
Coverage persists: if some $I\subseteq S$ had $a_I>1$, that term alone is
now $1$; otherwise nothing changed for $S$.
\item \emph{Cited.} The companion expectation-threshold theorem
\cite[Theorem~1.1]{OpenAIThreshold26} (called the ``constant-factor
rounding theorem'' in the source's introduction) then gives an integral
cover of cost at most $1/2$ at density
$\sigma/(25\cdot512^4)=p/(50\cdot512^4)$. Thus $C=50\cdot512^4$ suffices.
\qedhere
\end{enumerate}
\end{proof}

\unverified{Steps 2 and 4 quote \cite{Li26}, \cite{Park26} and
\cite{OpenAIThreshold26} exactly as the source uses them. The source
states neither their hypotheses nor the origin of $\sigma=p/2$, and I
could not check them against the cited papers. Steps 1 and 3 and the
arithmetic are verified.}

\section{An elementary covering lemma}
\label{sec:covering}

The lemma is applied in Section~\ref{sec:weights}, once per bin, with
$X=[N]$. Covers and costs are defined as in
Section~\ref{sec:introduction} with $X$ in place of $[N]$. An
\emph{edge} is a member of $\mathcal H$.

\begin{lemma}\label{lem:covering}
Let \(X\) be a finite set, let \(0<\rho<1\), and let \(r,h\) be integers
with \(1\le r\le h\) and \(h\ge64\). Suppose that
\(\mathcal H\subseteq\binom{X}{r}\) satisfies
\[
|\mathcal H|\rho^r\le2^{h+1}.
\]
Then the family
\[
\left\{S\subseteq X:
  |\{e\in\mathcal H:e\subseteq S\}|\ge2^{12h}\right\}
\]
has a cover \(\mathcal G\subseteq2^X\) with
\[
c_\rho(\mathcal G)\le3\cdot2^{-h}.
\]
\end{lemma}

\begin{proof}
\emph{Plan.} We build
$\mathcal G=\mathcal C_0\cup\mathcal C_1\cup\mathcal C_2$, where
$\mathcal C_0$ is deterministic and $\mathcal C_1,\mathcal C_2$ depend on
a random sample. We show that $\mathcal G$ covers the target family for
every outcome (Step~4) and that
$\mathbb E[c_\rho(\mathcal G)]\le3\cdot2^{-h}$ (Steps 1, 3, 6).
Throughout,
\[
m=2^{12h},\qquad T=2^{4h},\qquad \pi=2^{-8h}.
\]

\step{Step 1: sets of large weighted degree.}
For $I\subseteq X$ with $|I|\le r$ put
\[
d(I)=|\{e\in\mathcal H:I\subseteq e\}|\,\rho^{r-|I|},
\qquad
\mathcal C_0=\{I\subseteq X:|I|\le r,\ d(I)>T\}.
\]
Since $d(I)\rho^{|I|}=|\{e\in\mathcal H:I\subseteq e\}|\rho^r$ and each
edge has $2^r$ subsets,
$\sum_{|I|\le r}d(I)\rho^{|I|}=2^r|\mathcal H|\rho^r$. Each
$I\in\mathcal C_0$ satisfies $\rho^{|I|}<T^{-1}d(I)\rho^{|I|}$, so, by
the hypothesis on $|\mathcal H|\rho^r$ and $r\le h$,
\begin{equation}\label{eq:covering-initial-cost}
c_\rho(\mathcal C_0)
\le T^{-1}2^r|\mathcal H|\rho^r
\le2^{-4h}\,2^r\,2^{h+1}=2^{r+1-3h}\le2^{1-2h}.
\end{equation}

\step{Step 2: regular edges and the extension bound.}
Call an edge \emph{regular} if it contains no member of $\mathcal C_0$,
and let $\mathcal H'$ be the set of regular edges. Thus $d(I)\le T$
whenever $I$ is a subset of a regular edge. We claim that, for every
$B\subseteq X$,
\begin{equation}\label{eq:covering-extension}
\sum_{e\in\mathcal H'}\rho^{|e\setminus B|}
\le T\sum_{i=0}^r\binom{|B|}{i}.
\end{equation}
Group the edges $e\in\mathcal H'$ by $I=e\cap B$, a subset of $B$ of
size at most $r$; there are at most $\sum_{i=0}^r\binom{|B|}{i}$ groups.
Every edge in the group of $I$ has $|e\setminus B|=r-|I|$, so a nonempty
group contributes
\[
|\{e\in\mathcal H':e\cap B=I\}|\,\rho^{r-|I|}\le d(I)\le T,
\]
the last step because $I$ lies in a regular edge.

\step{Step 3: sampled pairs.}
Independently for each ordered pair $(e,e')\in(\mathcal H')^2$, diagonal
pairs included, \emph{sample} the pair with probability $\pi$. Let
$\mathcal C_1$ be the set of unions $e\cup e'$ of sampled pairs.
By \eqref{eq:covering-extension} with $B=e$ (so $|B|=r$ and the binomial
sum is $2^r$), $\sum_{e'\in\mathcal H'}\rho^{|e'\setminus e|}\le2^rT$.
Using $|e\cup e'|=r+|e'\setminus e|$ and linearity of expectation (a
union arising from several pairs is counted several times, which only
increases the bound),
\begin{equation}\label{eq:covering-pair-cost}
\mathbb E[c_\rho(\mathcal C_1)]
\le\pi\sum_{e\in\mathcal H'}\rho^r\sum_{e'\in\mathcal H'}\rho^{|e'\setminus e|}
\le\pi|\mathcal H'|\rho^r2^rT
\le2^{-8h}\,2^{h+1}\,2^r\,2^{4h}\le2^{1-2h}.
\end{equation}

\step{Step 4: residual generators; $\mathcal G$ is a cover.}
For every ordered $m$-tuple $(e_1,\ldots,e_m)$ of distinct regular edges
such that none of the pairs $(e_i,e_j)$, $1\le i,j\le m$, was sampled,
put $e_1\cup\cdots\cup e_m$ into $\mathcal C_2$. Now fix any outcome and
let $S$ contain at least $m$ edges but no member of $\mathcal C_0$.
Every edge $e\subseteq S$ is regular, since a member of $\mathcal C_0$
inside $e$ would lie inside $S$. Pick $m$ distinct edges
$e_1,\ldots,e_m\subseteq S$. If some pair $(e_i,e_j)$ was sampled, then
$e_i\cup e_j\in\mathcal C_1$ lies in $S$; otherwise
$e_1\cup\cdots\cup e_m\in\mathcal C_2$ lies in $S$.

\step{Step 5: total weight of tuples.}
For $0\le t\le m$ let
\[
W_t=
\sum_{\substack{(e_1,\ldots,e_t)\in(\mathcal H')^t\\ \text{distinct}}}
  \rho^{|e_1\cup\cdots\cup e_t|},
\qquad W_0=1.
\]
We claim $W_{t+1}\le T(1+rm)^rW_t$ for $t<m$. Fix a $t$-tuple with union
$B$; then $|B|\le rt\le rm$. Its extensions are indexed by regular edges
$e$ distinct from $e_1,\ldots,e_t$; summing over all regular $e$ only
increases the sum. Since
$\rho^{|B\cup e|}=\rho^{|B|}\rho^{|e\setminus B|}$,
\eqref{eq:covering-extension} gives
\[
\sum_{e\in\mathcal H'}\rho^{|B\cup e|}
\le\rho^{|B|}\,T\sum_{i=0}^r\binom{|B|}{i}
\le\rho^{|B|}\,T(1+rm)^r.
\]
The last step is $\sum_{i=0}^r\binom ni\le(1+n)^r$: list a subset of size
at most $r$ of an ordered $n$-set in increasing order and pad it to
length $r$ with a new symbol; this is an injection into words of length
$r$ over $n+1$ symbols. Summing over $t$-tuples proves the claim, and
iterating gives $W_m\le[T(1+rm)^r]^m$.

\step{Step 6: expected cost of $\mathcal C_2$.}
Fix an $m$-tuple of distinct regular edges. Its $m^2$ ordered pairs
$(e_i,e_j)$ are pairwise distinct, hence sampled independently, so the
tuple contributes to $\mathcal C_2$ with probability
$(1-\pi)^{m^2}\le e^{-\pi m^2}\le2^{-\pi m^2}$. (This is the only use of
independence.) By linearity of expectation, again counting with
multiplicity,
\[
\mathbb E[c_\rho(\mathcal C_2)]
\le W_m\,2^{-\pi m^2}\le[T(1+rm)^r]^m\,2^{-\pi m^2}.
\]
Now $1+rm\le(h+1)2^{12h}\le2^{13h}$ by $r\le h$ and $h+1\le2^h$, so
$T(1+rm)^r\le2^{4h+13rh}$. Also $\pi m=2^{4h}$ and
\[
4h+13rh+1\le5h^2+13h^2=18h^2\le18\cdot2^{2h}\le2^{4h},
\]
using $4h+1\le5h^2$, $r\le h$, $h\le2^h$, and $18\le2^{2h}$ (as
$h\ge64$). Therefore
\begin{equation}\label{eq:covering-residual-cost}
\mathbb E[c_\rho(\mathcal C_2)]
\le2^{m(4h+13rh-\pi m)}\le2^{-m}.
\end{equation}

\step{Step 7: conclusion.}
By \eqref{eq:covering-initial-cost}, \eqref{eq:covering-pair-cost} and
\eqref{eq:covering-residual-cost},
$\mathbb E[c_\rho(\mathcal G)]\le2^{2-2h}+2^{-m}\le2\cdot2^{-h}
\le3\cdot2^{-h}$, since $h\ge2$ and $m\ge h$. The sample space is finite,
so some outcome has cost at most the expectation, and by Step~4 its
$\mathcal G$ is a cover.
\end{proof}

\section{Weights for exceptional sets}
\label{sec:weights}

The next proposition gives a cover at a fixed fraction of the density.
It is used in Section~\ref{sec:density} with $q=Lp$.

\begin{proposition}\label{prop:half-density}
Let $L=2^{70}$. For every integer $N\ge1$, every $q\in(0,1)$, and every
$\mathcal F\subseteq 2^{[N]}$ satisfying $\mu_q(\mathcal F)\ge1/2$,
the family $E_{32}(\mathcal F)$ is $(q/L)$-small.
\end{proposition}

\begin{proof}
\emph{Plan.} Define weights $w(U)$ (Step 1); prove a global bound
(Step 2) and a local bound on every exceptional set (Step 3); then build
the cover from one family of heavy sets and one application of
Lemma~\ref{lem:covering} per bin (Steps 4--7).

\step{Step 1: weights.}
Identify subsets of $[N]$ with indicator vectors and write
$f=\mathbf 1_{\mathcal F}$. For $U\subseteq[N]$ define
\begin{equation}\label{eq:weights-definition}
 b(U)=(1-q)^{|U|}
 \sum_{z\in\{0,1\}^{U}}(-1)^{|z|}\,
 \mathbb E[f(z,Y)],
 \qquad w(U)=b(U)^2,
\end{equation}
where $|z|=\sum_{i\in U}z_i$, the vector $Y$ has independent
Bernoulli-$q$ coordinates on $[N]\setminus U$, and $(z,Y)$ equals $z$ on
$U$ and $Y$ off $U$. In particular
$b(\varnothing)=\mu_q(\mathcal F)\ge1/2$.

\step{Step 2: global bound.} We claim
\begin{equation}\label{eq:global-weight-bound}
 \sum_{U\subseteq[N]}q^{|U|}w(U)\le1.
\end{equation}
Let $X\sim\mu_q$. The functions
\[
 \chi_U(x)=\prod_{i\in U}\frac{q-x_i}{\sqrt{q(1-q)}},
 \qquad U\subseteq[N],
\]
form an orthonormal basis of the real functions on $\{0,1\}^N$ under
$\mu_q$ (the biased Fourier basis \cite[Section~8.4]{ODonnell14}): each
factor has mean $0$ and variance $1$, the coordinates are independent,
and there are $2^N$ functions. With
$\widehat f(U)=\mathbb E[f(X)\chi_U(X)]$, Parseval's identity gives
$\sum_U\widehat f(U)^2=\mathbb E[f(X)^2]=\mu_q(\mathcal F)\le1$.
To relate $\widehat f$ to $b$, condition on the restriction of $X$ to $U$:
\[
 \widehat f(U)=\sum_{z\in\{0,1\}^U}
 \Bigl(\prod_{i\in U}\frac{\mathbb P(X_i=z_i)(q-z_i)}{\sqrt{q(1-q)}}\Bigr)
 \mathbb E[f(z,Y)].
\]
Both cases $z_i=0,1$ give
$\mathbb P(X_i=z_i)(q-z_i)=q(1-q)(-1)^{z_i}$ (namely $(1-q)q$ and
$q(q-1)$). Hence
\[
 \widehat f(U)
 =[q(1-q)]^{|U|/2}
 \sum_{z\in\{0,1\}^{U}}(-1)^{|z|}\,\mathbb E[f(z,Y)]
 =\left(\frac{q}{1-q}\right)^{|U|/2}b(U).
\]
Squaring, $q^{|U|}w(U)=(1-q)^{|U|}\widehat f(U)^2\le\widehat f(U)^2$;
summing over $U$ gives \eqref{eq:global-weight-bound}.

\step{Step 3: local bound.}
Fix $S\in E_{32}(\mathcal F)$. We claim
\begin{equation}\label{eq:exceptional-weight-bound}
 \sum_{\varnothing\ne U\subseteq S}w(U)^{16}\ge2^{-32}.
\end{equation}
Consider arrays $x\in\{0,1\}^{32\times N}$ with rows
$x_1,\ldots,x_{32}\in\{0,1\}^N$ (subsets of $[N]$) and columns
$x^{(1)},\ldots,x^{(N)}\in\{0,1\}^{32}$. For $y\in\{0,1\}^{32}$ with
$|y|=\sum_jy_j$, set
\[
 P(y)=q^{|y|}(1-q)^{32-|y|},
 \qquad
 Q(y)=P(y)-(1-q)^{32}(-1)^{|y|},
 \qquad
 \nu_S(x)=\prod_{i\notin S}P(x^{(i)})\prod_{i\in S}Q(x^{(i)}).
\]
($Q$ is the source's $R$; $R$ denotes random sets in
Section~\ref{sec:density}.)

\emph{(a) Vanishing.} We have
$\sum_x\nu_S(x)\prod_{j=1}^{32}f(x_j)=0$, term by term. Indeed
$Q(0,\ldots,0)=0$, so if $\nu_S(x)\ne0$ then every column in $S$ contains
a $1$, that is, $S\subseteq x_1\cup\cdots\cup x_{32}$. Since
$S\in E_{32}(\mathcal F)$, some row is not in $\mathcal F$, and
$\prod_jf(x_j)=0$. (The signs of $\nu_S$ play no role.)

\emph{(b) Expansion.} Expand $\prod_{i\in S}Q(x^{(i)})$ by choosing, in
each column $i\in S$, either $P(x^{(i)})$ or
$-(1-q)^{32}(-1)^{|x^{(i)}|}$; let $U\subseteq S$ be the set of columns
where the second term is chosen. Then
\[
 \nu_S(x)=\sum_{U\subseteq S}(-1)^{|U|}(1-q)^{32|U|}
 \prod_{i\in U}(-1)^{|x^{(i)}|}\prod_{i\notin U}P(x^{(i)}).
\]
Each factor splits over the $32$ rows:
$(1-q)^{32|U|}=\prod_j(1-q)^{|U|}$,
$\prod_{i\in U}(-1)^{|x^{(i)}|}=\prod_j(-1)^{\sum_{i\in U}x_{ji}}$, and
$\prod_{i\notin U}P(x^{(i)})
=\prod_j\prod_{i\notin U}q^{x_{ji}}(1-q)^{1-x_{ji}}$.
So $\nu_S(x)=\sum_{U\subseteq S}(-1)^{|U|}\prod_{j=1}^{32}g_U(x_j)$ with
\[
 g_U(u)=(1-q)^{|U|}(-1)^{\sum_{i\in U}u_i}
 \prod_{i\notin U}q^{u_i}(1-q)^{1-u_i}
 \qquad(u\in\{0,1\}^N).
\]
The rows of $x$ range independently over $\{0,1\}^N$, so
$\sum_x\prod_jg_U(x_j)f(x_j)=\bigl(\sum_ug_U(u)f(u)\bigr)^{32}$. Writing
$u=(z,y)$ with $z$ on $U$ and $y$ off $U$, the sum over $y$ against the
Bernoulli masses is $\mathbb E[f(z,Y)]$, so $\sum_ug_U(u)f(u)=b(U)$ by
\eqref{eq:weights-definition}. With (a),
\begin{equation}\label{eq:signed-identity}
 0=\sum_{U\subseteq S}(-1)^{|U|}b(U)^{32}.
\end{equation}

\emph{(c) Conclusion.} Isolating $U=\varnothing$ and using the triangle
inequality,
\[
 2^{-32}\le b(\varnothing)^{32}
 \le\sum_{\varnothing\ne U\subseteq S}|b(U)|^{32}
 =\sum_{\varnothing\ne U\subseteq S}w(U)^{16}.
\]

\step{Step 4: heavy sets.}
Set $\rho=q/L$ and
\[
 \mathcal G_0=
 \{U\subseteq[N]:U\ne\varnothing,\ w(U)\ge2^{-64|U|}\}.
\]
For $U\in\mathcal G_0$ with $|U|=r\ge1$,
$\rho^r=2^{-6r}q^r2^{-64r}\le2^{-6r}q^rw(U)\le2^{-6}q^rw(U)$, so by
\eqref{eq:global-weight-bound}
\begin{equation}\label{eq:outright-cost}
 c_\rho(\mathcal G_0)\le2^{-6}\sum_Uq^{|U|}w(U)\le2^{-6}.
\end{equation}

\step{Step 5: bins.}
For integers $r\ge1$ and $h$ (of any sign), let
\[
 \mathcal H_{r,h}=
 \bigl\{U\subseteq[N]:|U|=r,\ U\notin\mathcal G_0,
 \ 2^{-(h+1)}<w(U)\le2^{-h}\bigr\}.
\]
Every nonempty $U\notin\mathcal G_0$ with $w(U)>0$ lies in exactly one
bin. A nonempty bin has $h\ge64r$: a member satisfies
$2^{-(h+1)}<w(U)<2^{-64r}$, so $h+1>64r$. By
\eqref{eq:global-weight-bound},
\[
 |\mathcal H_{r,h}|\,q^r2^{-(h+1)}
 <\sum_{U\in\mathcal H_{r,h}}q^r w(U)\le1,
\]
and $\rho^r\le q^r$ because $L\ge1$; hence
$|\mathcal H_{r,h}|\rho^r\le2^{h+1}$. So Lemma~\ref{lem:covering} applies
with $X=[N]$ and $\mathcal H=\mathcal H_{r,h}$; its hypotheses
$1\le r\le h$ and $h\ge64$ follow from $h\ge64r\ge64$. It gives a family
$\mathcal G_{r,h}$ with $c_\rho(\mathcal G_{r,h})\le3\cdot2^{-h}$ that
covers every set containing at least $2^{12h}$ members of
$\mathcal H_{r,h}$.

\step{Step 6: the cover and its cost.}
Let $\mathcal G$ be the union of $\mathcal G_0$ and the
$\mathcal G_{r,h}$ over the nonempty bins (finitely many, since $[N]$ is
finite). As each nonempty bin has $h\ge64r$,
\begin{equation}\label{eq:bin-cost-sum}
 c_\rho(\mathcal G)
 \le2^{-6}+3\sum_{r\ge1}\sum_{h\ge64r}2^{-h}
 =\frac1{64}+3\sum_{r\ge1}2^{1-64r}
 =\frac1{64}+\frac{6}{2^{64}-1}<\frac12.
\end{equation}

\step{Step 7: $\mathcal G$ covers $E_{32}(\mathcal F)$.}
Suppose $S\subseteq[N]$ contains no member of $\mathcal G$. Then $S$
contains no member of $\mathcal G_0$ and, for each nonempty bin, fewer
than $2^{12h}$ members of $\mathcal H_{r,h}$ (otherwise it would contain
a member of $\mathcal G_{r,h}$). So every nonempty $U\subseteq S$ has
$w(U)=0$ or lies in a bin, where $w(U)^{16}\le2^{-16h}$. Therefore
\begin{equation}\label{eq:bin-power-sum}
 \sum_{\varnothing\ne U\subseteq S}w(U)^{16}
 \le\sum_{r\ge1}\sum_{h\ge64r}2^{12h}\cdot2^{-16h}
 =\sum_{r\ge1}\frac{16}{15}\,2^{-256r}
 =\frac{16}{15(2^{256}-1)}<2^{-32}.
\end{equation}
By \eqref{eq:exceptional-weight-bound}, $S\notin E_{32}(\mathcal F)$.
Contrapositively, every member of $E_{32}(\mathcal F)$ contains a member
of $\mathcal G$; with \eqref{eq:bin-cost-sum}, $E_{32}(\mathcal F)$ is
$\rho$-small.
\end{proof}

\section{Returning to the original density}
\label{sec:density}

\begin{proof}[Proof of Theorem~\ref{thm:main}]
\emph{Plan.} Build $L=2^{70}$ dependent random sets, each with law
exactly $\mu_p$, whose union has density $Lp$ (or is all of $[N]$). A
union bound puts all of them in $\mathcal D$ simultaneously with high
probability. Then apply Proposition~\ref{prop:half-density} to unions of
$L$ members of $\mathcal D$. (The source's Case~2 parameter $t$ is
called $\theta$ here.)

Let $L=2^{70}$, so $k=32L$, and assume $\mu_p(\mathcal D)\ge1-1/k$. If
$R_1,\ldots,R_L$ are random subsets of $[N]$ on one probability space,
each with law $\mu_p$, then, with no independence assumption,
\begin{equation}
\label{eq:all-good-rows}
\mathbb P(R_j\in\mathcal D\text{ for all }1\le j\le L)
\ge1-\sum_{j=1}^L\mathbb P(R_j\notin\mathcal D)
\ge1-\frac Lk=\frac{31}{32}.
\end{equation}

\medskip
\noindent\emph{Case 1: $Lp<1$.}
Independently for each $i\in[N]$, choose a label
$J_i\in\{0,1,\ldots,L\}$ with
\[
\mathbb P(J_i=0)=1-Lp,
\qquad
\mathbb P(J_i=j)=p\quad(1\le j\le L),
\]
and put $R_j=\{i:J_i=j\}$, so each element goes to at most one row. For
fixed $j$, the events $\{i\in R_j\}$ are independent over $i$, each of
probability $p$; so $R_j$ has law $\mu_p$. The union
$R=R_1\cup\cdots\cup R_L=\{i:J_i\ne0\}$ has independent coordinates,
each present with probability $q:=Lp\in(0,1)$; so $R$ has law $\mu_q$.

Define
$\mathcal F=\{D_1\cup\cdots\cup D_L:D_1,\ldots,D_L\in\mathcal D\}$.
On the event in~\eqref{eq:all-good-rows}, $R\in\mathcal F$, hence
\[
\mu_q(\mathcal F)=\mathbb P(R\in\mathcal F)\ge\frac{31}{32}\ge\frac12.
\]
The unions of $32$ members of $\mathcal F$ are exactly the unions of
$32L$ members of $\mathcal D$: expand each member of $\mathcal F$, or
conversely group a $32L$-tuple into $32$ blocks of length $L$. Thus
\[
E_{32}(\mathcal F)=E_{32L}(\mathcal D)=E_k(\mathcal D).
\]
Proposition~\ref{prop:half-density} shows that this family is
$(q/L)$-small, and $q/L=p$.

\medskip
\noindent\emph{Case 2: $Lp\ge1$.}
Set
\[
\theta=\frac{p-1/L}{1-1/L},
\]
so $\theta\in[0,1)$ because $1/L\le p<1$. Independently for each
$i\in[N]$: choose one row uniformly among the $L$ rows and put $i$ in
it; then put $i$ into each other row independently with probability
$\theta$. For fixed $j$,
\[
\mathbb P(i\in R_j)=\frac1L+\left(1-\frac1L\right)\theta=p,
\]
and these events are independent over $i$, so $R_j$ has law $\mu_p$;
moreover $R_1\cup\cdots\cup R_L=[N]$ in every outcome.
By~\eqref{eq:all-good-rows}, some outcome has all $R_j\in\mathcal D$.
Repeating these $L$ sets $32$ times gives $k$ members of $\mathcal D$
whose union $[N]$ contains every $S$. Hence $E_k(\mathcal D)=\varnothing$,
which the empty family covers at cost $0$. (At the boundary $Lp=1$,
$\theta=0$.)
\end{proof}

\section{Verification checklist}

\begin{enumerate}
\item \emph{Extension bound} \eqref{eq:covering-extension}
(Lemma~\ref{lem:covering}, Step 2). Regularity is used exactly to get
$d(e\cap B)\le T$. This bound drives both
\eqref{eq:covering-pair-cost} and the recursion in Step 5.
\item \emph{Residual cost} (Lemma~\ref{lem:covering}, Step 6). The $m^2$
pairs of a tuple are distinct, giving probability $(1-\pi)^{m^2}$; then
the arithmetic $4h+13rh+1\le2^{4h}=\pi m$.
\item \emph{Coefficient identity}
(Proposition~\ref{prop:half-density}, Step 2).
$\widehat f(U)=(q/(1-q))^{|U|/2}b(U)$, and the discarded factor
$(1-q)^{|U|}\le1$ goes in the right direction.
\item \emph{Signed identity} \eqref{eq:signed-identity} (Step 3).
$Q(0,\ldots,0)=0$ gives termwise vanishing; the row factorisation
produces $b(U)^{32}$; $b(\varnothing)\ge1/2$.
\item \emph{Bins} (Steps 5--7). Nonempty bins have $h\ge64r$; the sums
\eqref{eq:bin-cost-sum} and \eqref{eq:bin-power-sum}; the threshold
$2^{12h}$ is the same in Lemma~\ref{lem:covering} and in Step 7.
\item \emph{Coupling} (Section~\ref{sec:density}). Each row has law
exactly $\mu_p$; the union has law $\mu_{Lp}$ (Case 1) or equals $[N]$
(Case 2); $E_{32}(\mathcal F)=E_k(\mathcal D)$ and $q/L=p$.
\item \emph{Corollary~\ref{cor:two-union} only.} It depends on cited
statements that are not checked here (see the boxed note);
Theorem~\ref{thm:main} does not use it.
\end{enumerate}

\begin{thebibliography}{99}

\bibitem{AscoliEtAl26}
Ruben Ascoli, Xiaoyu He, Jinyoung Park, and Michel Talagrand,
\emph{A reformulation of the discrete Convexity Conjecture via
$k$-thresholds},
arXiv:2608.11183v1, 2026.
\href{https://arxiv.org/abs/2608.11183v1}{arXiv:2608.11183v1}.

\bibitem{EllisFilmusFriedgut12}
David Ellis, Yuval Filmus, and Ehud Friedgut,
\emph{Triangle-intersecting families of graphs},
Journal of the European Mathematical Society \textbf{14} (2012),
no.~3, 841--885.
\href{https://doi.org/10.4171/JEMS/320}{doi:10.4171/JEMS/320}.

\bibitem{FischerPerson25}
Thomas Fischer and Yury Person,
\emph{Further remarks on fractional vs. expectation thresholds},
arXiv:2505.21782v1, 2025.
\href{https://doi.org/10.48550/arXiv.2505.21782}{doi:10.48550/arXiv.2505.21782}.

\bibitem{Friedgut08}
Ehud Friedgut,
\emph{On the measure of intersecting families, uniqueness and stability},
Combinatorica \textbf{28} (2008), no.~5, 503--528.
\href{https://doi.org/10.1007/s00493-008-2318-9}{doi:10.1007/s00493-008-2318-9}.

\bibitem{HuaSongTudose26}
Dongming (Merrick) Hua, Antoine Song, and Stefan Tudose,
\emph{On Talagrand's Convexity Conjecture},
arXiv:2605.10908v2, 2026.
\href{https://arxiv.org/abs/2605.10908v2}{arXiv:2605.10908v2}.

\bibitem{Li26}
Chen Li,
\emph{On $p$-spread measures},
arXiv:2609.08967v1, 2026.
\href{https://arxiv.org/abs/2609.08967v1}{arXiv:2609.08967v1}.

\bibitem{ODonnell14}
Ryan O'Donnell,
\emph{Analysis of Boolean Functions},
Cambridge University Press, Cambridge, 2014.
Author's corrected text, 2021:
\href{https://arxiv.org/abs/2105.10386v1}{arXiv:2105.10386v1}.

\bibitem{OpenAIThreshold26}
OpenAI,
\emph{Integral and fractional expectation thresholds are equivalent},
OpenAI Math Release preprint
\href{https://github.com/openai/math/blob/main/preprints/Integral-and-fractional-expectation-thresholds-are-equivalent-September-23-2026/paper.pdf}{OAI:Integral-and-fractional-expectation-thresholds-are-equivalent-September-23-2026},
2026.

\bibitem{Park26}
Jinyoung Park,
\emph{A dimension-free comparison between expectation thresholds and
fractional expectation thresholds},
arXiv:2609.14681v1, 2026.
\href{https://arxiv.org/abs/2609.14681v1}{arXiv:2609.14681v1}.

\bibitem{ParkPham24}
Jinyoung Park and Huy Tuan Pham,
\emph{On a conjecture of Talagrand on selector processes and a consequence
on positive empirical processes},
Annals of Mathematics \textbf{199} (2024), no.~3, 1293--1321.
\href{https://doi.org/10.4007/annals.2024.199.3.6}{doi:10.4007/annals.2024.199.3.6}.

\bibitem{Talagrand10}
Michel Talagrand,
\emph{Are many small sets explicitly small?},
Proceedings of the Forty-Second ACM Symposium on Theory of Computing,
ACM, New York, 2010, 13--36.
\href{https://doi.org/10.1145/1806689.1806693}{doi:10.1145/1806689.1806693}.

\bibitem{Talagrand21}
Michel Talagrand,
\emph{Upper and Lower Bounds for Stochastic Processes:
Decomposition Theorems},
2nd ed., Ergebnisse der Mathematik und ihrer Grenzgebiete, 3.~Folge,
vol.~60, Springer, Cham, 2021.
\href{https://doi.org/10.1007/978-3-030-82595-9}{doi:10.1007/978-3-030-82595-9}.

\end{thebibliography}

\end{document}
